\documentclass[10pt]{extarticle}
\usepackage[letterpaper, margin=1in]{geometry}
\usepackage{graphicx}
\usepackage{amsmath}
\usepackage{amssymb}
\usepackage{fancyhdr}
\usepackage{lastpage}
\usepackage{enumitem}
\usepackage{cancel}
\usepackage{tcolorbox}
\usepackage{bold-extra}
\usepackage{multicol}
\allowdisplaybreaks
\title{\centering\Large\textsc{\textbf{Applications of Linear Algebra: Problem Set 1}}}
\author{\large \textsc{Rento Saijo}}
\date{Feruary 4, 2026}
\pagestyle{fancy}
\fancyhf{}
\fancyfoot[C]{\textsc{Page \thepage\ of \pageref{LastPage}}}
\renewcommand{\headrulewidth}{0pt}
\tcbset{colframe=darkgray!100, arc=0.1mm, boxrule=2pt, boxsep=1mm}
\newtcolorbox{problem}[1][]{title={\textsc{#1}}}
\begin{document}
\maketitle
\thispagestyle{fancy}
\vspace{-1em}\noindent\rule{\linewidth}{0.6pt}\vspace{0.5em}

\begin{problem}
[Problem 1]
The following table shows the team home runs (HR) and team wins for the 2019 season for six teams in the American League of Major League Baseball: New York Yankees (NYY), Baltimore Orlioles (Balt), Cleveland Indians (Cleve), Chicago White Sox (CWS), Texas Rangers (TX), and Los Angeles Angels (LAA).
\begin{center}
\begin{tabular}{l|l|l|l|l|l|l}
Team & NYY & Balt & Cleve & CWS & TX & LAA \\ \hline
HR & 306 & 213 & 223 & 182 & 223 & 220 \\ \hline
Wins & 103 & 54 & 93 & 72 & 78 & 72
\end{tabular}
\end{center}
\begin{enumerate}[label=(\alph*)]
\item Before calculating anything, make a guess about the correlation coefficient between home runs and wins.
\item Now calculate the correlation coeffiecient between the two quantities. Did the result agree with your guess?
\end{enumerate}
\end{problem}

\begin{equation}
2+3 = 4
\label{cat}
\end{equation}

Equation~\ref{cat}

\begin{enumerate}[label=(\alph*)]
\item We'd expect a positive, moderately strong correlation: teams that hit more home runs tend to win more games (i.e., we'd expect a team with above average number of home runs to have above average number of wins). Perhaps, something along the lines of $r=0.75$ would be reasonable.
\item Let
\[
\vec{x}=\begin{bmatrix}306\\213\\223\\182\\223\\220\end{bmatrix},\qquad
\vec{y}=\begin{bmatrix}103\\54\\93\\72\\78\\72\end{bmatrix}.
\]
The Pearson correlation coefficient is
\[
r=\frac{(\vec{x}-\bar x\,\vec{\mathbf{1}})\cdot (\vec{y}-\bar y\,\vec{\mathbf{1}})}{\|\vec{x}-\bar x\,\vec{\mathbf{1}}\|\,\|\vec{y}-\bar y\,\vec{\mathbf{1}}\|}.
\]
Let us compute the means:
\[
\bar x=\frac{306+213+223+182+223+220}{6}=\frac{1367}{6},
\qquad
\bar y=\frac{103+54+93+72+78+72}{6}=\frac{472}{6}.
\]
Thus the de-meaned vectors are
\[
\vec{x}-\bar x\,\vec{\mathbf{1}}
=
\frac{1}{6}
\begin{bmatrix}
469\\
-89\\
-29\\
-275\\
-29\\
-47
\end{bmatrix},
\qquad
\vec{y}-\bar y\,\vec{\mathbf{1}}
=
\frac{1}{6}
\begin{bmatrix}
146\\
-236\\
86\\
-40\\
-4\\
-40
\end{bmatrix}.
\]
Let us compute the dot product:
\begin{align}
(\vec{x}-\bar x\,\vec{\mathbf{1}})\cdot (\vec{y}-\bar y\,\vec{\mathbf{1}})
&=\frac{1}{36}\Big(469\cdot146+(-89)(-236)+(-29)(86)+(-275)(-40)+(-29)(-4)+(-47)(-40)\Big) \notag \\
&=\frac{92148}{36}.
\end{align}
Let us compute the norms:
\[
\|\vec{x}-\bar x\,\vec{\mathbf{1}}\|
=
\sqrt{\frac{1}{36}\Big(469^2+89^2+29^2+275^2+29^2+47^2\Big)}
=
\sqrt{\frac{351390}{36}},
\]
\[
\|\vec{y}-\bar y\,\vec{\mathbf{1}}\|
=
\sqrt{\frac{1}{36}\Big(146^2+236^2+86^2+40^2+4^2+40^2\Big)}
=
\sqrt{\frac{47064}{36}}.
\]
Therefore,
\[
r
=
\frac{\frac{92148}{36}}{\sqrt{\frac{351390}{36}}\sqrt{\frac{47064}{36}}}
=
\frac{92148}{\sqrt{351390\cdot47064}}=0.716.
\]
So the correlation between home runs and wins is $\boxed{r=0.716}$, which is a fairly strong positive correlation and is consistent with the guess in part (a).
\end{enumerate}

\begin{problem}
[Problem 2]
The table here shows a term-document array based on a library of five mathematical books.
\begin{minipage}{0.45\textwidth}
\hfill
\begin{tabular}{l | l | l | l | l | l}
& B1 & B2 & B3 & B4 & B5 \\ \hline
Dimension & 0 & 1 & 1 & 0 & 0\\ \hline
Geometry & 1 & 1 & 1 & 0 & 1\\ \hline
Probability & 1 & 1 & 1 & 1 & 0 \\ \hline
Statistics & 0 & 1 & 0 & 1 & 0 \\ \hline
Symmetry & 0 & 0 & 1 & 0 & 1 \\ \hline
\end{tabular}
\end{minipage}
\begin{minipage}{0.5\textwidth}
\begin{itemize}
\item[B1] Dunham, \emph{Euler: The Master of Us All};
\item[B2] Osserman, \emph{Poetry of the Unvierse};
\item[B3] Pappas, \emph{The Joy of Mathematics};
\item[B4] P\"olya, \emph{Patterns of Plausible Inference};
\item[B5] Yaglom, \emph{Geometric Transformations I}
\end{itemize}
\end{minipage}
\medskip

For each query $\mathbf{q}$ given below, apply the algorithm from class to determine which documents are relevant to the query with a threshold of (a) $\alpha=0.5$ and (b) $\alpha=0.8$.
\begin{enumerate}[label=\roman*.]
\item $\mathbf{q}$ is ``geometry" and ``probability"
\item $\mathbf{q}$ is ``geometry" and ``symmetry"
\item $\mathbf{q}$ is ``dimension" and ``statistics"
\end{enumerate}
\end{problem}

\begin{enumerate}[label=\roman*.]
\item Fix the term order as
\[
(\text{Dimension},\text{Geometry},\text{Probability},\text{Statistics},\text{Symmetry}).
\]
Thus, the document vectors are
\[
\vec{a}_1=\begin{bmatrix}0\\1\\1\\0\\0\end{bmatrix},\quad
\vec{a}_2=\begin{bmatrix}1\\1\\1\\1\\0\end{bmatrix},\quad
\vec{a}_3=\begin{bmatrix}1\\1\\1\\0\\1\end{bmatrix},\quad
\vec{a}_4=\begin{bmatrix}0\\0\\1\\1\\0\end{bmatrix},\quad
\vec{a}_5=\begin{bmatrix}0\\1\\0\\0\\1\end{bmatrix}.
\]
We compute
\[
\cos\theta=\frac{\vec{a}\cdot\vec{q}}{\|\vec{a}\|\|\vec{q}\|},
\]
and include a document when $\cos\theta>\alpha$.
For this query,
\[
\vec{q}=\begin{bmatrix}0\\1\\1\\0\\0\end{bmatrix},\qquad \|\vec{q}\|=\sqrt{2}.
\]
Also,
\[
\|\vec{a}_1\|=\sqrt{2},\qquad \|\vec{a}_2\|=2,\qquad \|\vec{a}_3\|=2,\qquad \|\vec{a}_4\|=\sqrt{2},\qquad \|\vec{a}_5\|=\sqrt{2}.
\]
Let us compute the dot products:
\[
\vec{a}_1\cdot\vec{q}=2,\qquad
\vec{a}_2\cdot\vec{q}=2,\qquad
\vec{a}_3\cdot\vec{q}=2,\qquad
\vec{a}_4\cdot\vec{q}=1,\qquad
\vec{a}_5\cdot\vec{q}=1.
\]
Thus
\begin{align*}
\cos\theta_1&=\frac{2}{\sqrt{2}\sqrt{2}}=1,\\
\cos\theta_2&=\frac{2}{2\sqrt{2}}=\frac{1}{\sqrt{2}},\qquad
\cos\theta_3=\frac{2}{2\sqrt{2}}=\frac{1}{\sqrt{2}},\\
\cos\theta_4&=\frac{1}{\sqrt{2}\sqrt{2}}=\frac{1}{2},\qquad
\cos\theta_5=\frac{1}{\sqrt{2}\sqrt{2}}=\frac{1}{2}.
\end{align*}
For $\alpha=0.5$, we include $B1,B2,B3$ (note $B4$ and $B5$ have $\cos\theta=\frac{1}{2}$, so they are not included since we require $\cos\theta>\alpha$).
For $\alpha=0.8$, we only include $B1$.
\item Similarly, for this query,
\[
\vec{q}=\begin{bmatrix}0\\1\\0\\0\\1\end{bmatrix},\qquad \|\vec{q}\|=\sqrt{2}.
\]
Let us compute the dot products:
\[
\vec{a}_1\cdot\vec{q}=1,\qquad
\vec{a}_2\cdot\vec{q}=1,\qquad
\vec{a}_3\cdot\vec{q}=2,\qquad
\vec{a}_4\cdot\vec{q}=0,\qquad
\vec{a}_5\cdot\vec{q}=2.
\]
Thus
\[
\cos\theta_1=\frac{1}{\sqrt{2}\sqrt{2}}=\frac{1}{2},\qquad
\cos\theta_2=\frac{1}{2\sqrt{2}},\qquad
\cos\theta_3=\frac{2}{2\sqrt{2}}=\frac{1}{\sqrt{2}},\qquad
\cos\theta_4=0,\qquad
\cos\theta_5=\frac{2}{\sqrt{2}\sqrt{2}}=1.
\]
For $\alpha=0.5$, we include $B3,B5$ (note $B1$ has $\cos\theta=\frac{1}{2}$, so it is not included).
For $\alpha=0.8$, we only include $B5$.
\item Similarly, for this query,
\[
\vec{q}=\begin{bmatrix}1\\0\\0\\1\\0\end{bmatrix},\qquad \|\vec{q}\|=\sqrt{2}.
\]
Let us compute the dot products:
\[
\vec{a}_1\cdot\vec{q}=0,\qquad
\vec{a}_2\cdot\vec{q}=2,\qquad
\vec{a}_3\cdot\vec{q}=1,\qquad
\vec{a}_4\cdot\vec{q}=1,\qquad
\vec{a}_5\cdot\vec{q}=0.
\]
Thus
\[
\cos\theta_1=0,\qquad
\cos\theta_2=\frac{2}{2\sqrt{2}}=\frac{1}{\sqrt{2}},\qquad
\cos\theta_3=\frac{1}{2\sqrt{2}},\qquad
\cos\theta_4=\frac{1}{\sqrt{2}\sqrt{2}}=\frac{1}{2},\qquad
\cos\theta_5=0.
\]
For $\alpha=0.5$, we only include $B2$ (note $B4$ has $\cos\theta=\frac{1}{2}$, so it is not included).
For $\alpha=0.8$, we won't include any.
\end{enumerate}

\begin{problem}
[Problem 3]
Let $\vec{x}$ be a block vector with two vector elements, $\vec{x}=\begin{bmatrix}\vec{a}\\ \vec{b}\end{bmatrix}$, where $\vec{a}\in\mathbb{R}^N$ and $\vec{b}\in\mathbb{R}^M$.
\begin{enumerate}[label=(\alph*)]
\item Express $\mathrm{rms}(\vec{x})$ in terms of $\mathrm{rms}(\vec{a})$ and $\mathrm{rms}(\vec{b})$.
\item Express $\mathrm{avg}(\vec{x})$ in terms of $\mathrm{avg}(\vec{a})$ and $\mathrm{avg}(\vec{b})$.
\end{enumerate}
\end{problem}

\begin{enumerate}[label=(\alph*)]
\item Since $\vec{a}$ has $N$ entries and $\vec{b}$ has $M$ entries, the stacked vector $\vec{x}=\begin{bmatrix}\vec{a}\\ \vec{b}\end{bmatrix}$ has $N+M$ entries. By definition,
\[
\mathrm{rms}(\vec{x})=\frac{\|\vec{x}\|}{\sqrt{N+M}}=\sqrt{\frac{1}{N+M}\|\vec{x}\|^2}.
\]
Since $\|\vec{x}\|^2=\sum_{i=1}^{N+M}x_i^2$, we have
\[
\mathrm{rms}(\vec{x})=\sqrt{\frac{1}{N+M}\sum_{i=1}^{N+M}x_i^2}.
\]
Since $\vec{x}$ is $\vec{a}$ stacked on $\vec{b}$,
\[
\sum_{i=1}^{N+M}x_i^2=\sum_{i=1}^{N}a_i^2+\sum_{j=1}^{M}b_j^2.
\]
Again, by definition,
\[
\mathrm{rms}(\vec{a})=\sqrt{\frac{1}{N}\sum_{i=1}^{N}a_i^2}
\qquad\Rightarrow\qquad
\mathrm{rms}(\vec{a})^2=\frac{1}{N}\sum_{i=1}^{N}a_i^2
\qquad\Rightarrow\qquad
\sum_{i=1}^{N}a_i^2=N\,\mathrm{rms}(\vec{a})^2,
\]
\[
\mathrm{rms}(\vec{b})=\sqrt{\frac{1}{M}\sum_{j=1}^{M}b_j^2}
\qquad\Rightarrow\qquad
\mathrm{rms}(\vec{b})^2=\frac{1}{M}\sum_{j=1}^{M}b_j^2
\qquad\Rightarrow\qquad
\sum_{j=1}^{M}b_j^2=M\,\mathrm{rms}(\vec{b})^2.
\]
Substituting these into the expression
\[
\sum_{i=1}^{N+M}x_i^2=\sum_{i=1}^{N}a_i^2+\sum_{j=1}^{M}b_j^2
\]
gives
\[
\sum_{i=1}^{N+M}x_i^2=N\,\mathrm{rms}(\vec{a})^2+M\,\mathrm{rms}(\vec{b})^2.
\]
Substituting this into
\[
\mathrm{rms}(\vec{x})=\sqrt{\frac{1}{N+M}\sum_{i=1}^{N+M}x_i^2}
\]
gives
\[
\mathrm{rms}(\vec{x})=\sqrt{\frac{N\,\mathrm{rms}(\vec{a})^2+M\,\mathrm{rms}(\vec{b})^2}{N+M}}.
\]
\item Since $\vec{a}$ has $N$ entries and $\vec{b}$ has $M$ entries, the stacked vector $\vec{x}=\begin{bmatrix}\vec{a}\\ \vec{b}\end{bmatrix}$ has $N+M$ entries. By definition,
\[
\mathrm{avg}(\vec{x})=\frac{1}{N+M}\vec{\mathbf{1}}\cdot\vec{x}.
\]
Since $\vec{x}$ is $\vec{a}$ stacked on $\vec{b}$, we can "decompose" (?):
\[
\vec{\mathbf{1}}=\begin{bmatrix}\vec{\mathbf{1}}\\ \vec{\mathbf{1}}\end{bmatrix}.
\]
Thus,
\[
\vec{\mathbf{1}}\cdot\vec{x}
=
\begin{bmatrix}\vec{\mathbf{1}}\\ \vec{\mathbf{1}}\end{bmatrix}\cdot
\begin{bmatrix}\vec{a}\\ \vec{b}\end{bmatrix}
=
\vec{\mathbf{1}}\cdot\vec{a}+\vec{\mathbf{1}}\cdot\vec{b}.
\]
Thus,
\[
\mathrm{avg}(\vec{x})=\frac{1}{N+M}\left(\vec{\mathbf{1}}\cdot\vec{a}+\vec{\mathbf{1}}\cdot\vec{b}\right).
\]
Again, by definition,
\[
\mathrm{avg}(\vec{a})=\frac{1}{N}\vec{\mathbf{1}}\cdot\vec{a}
\qquad\Rightarrow\qquad
\vec{\mathbf{1}}\cdot\vec{a}=N\,\mathrm{avg}(\vec{a}),
\]
\[
\mathrm{avg}(\vec{b})=\frac{1}{M}\vec{\mathbf{1}}\cdot\vec{b}
\qquad\Rightarrow\qquad
\vec{\mathbf{1}}\cdot\vec{b}=M\,\mathrm{avg}(\vec{b}).
\]
Substituting these into
\[
\mathrm{avg}(\vec{x})=\frac{1}{N+M}\left(\vec{\mathbf{1}}\cdot\vec{a}+\vec{\mathbf{1}}\cdot\vec{b}\right)
\]
gives
\[
\mathrm{avg}(\vec{x})=\frac{N\,\mathrm{avg}(\vec{a})+M\,\mathrm{avg}(\vec{b})}{N+M}.
\]
\end{enumerate}

\begin{problem}
[Problem 4]
As mentioned in class, any real-valued function $f:\mathbb{R}^N\to\mathbb{R}$ that satisfies (i) nonnegative homogeneity, (ii) the triangle inequality, and (iii) nonnegative definiteness is called a \emph{vector norm}. The most common vector norm is the Euclidean norm which is sometimes called the two-norm, and is denoted $\|\vec{x}\|_2$. Show the following are also vector norms.
\begin{multicols}{2}
\begin{enumerate}[label=(\alph*)]
\item $\|\vec{x}\|_1=|x_1|+\cdots+|x_N|$
\item $\|\vec{x}\|_\infty=\max\{|x_1|,\dots,|x_N|\}$
\end{enumerate}
\end{multicols}
\end{problem}

\begin{enumerate}[label=(\alph*)]
\item Let us check the three axioms for the one-norm.
\begin{enumerate}[label=(\roman*)]
\item For any scalar $c\in\mathbb{R}$,
\[
\|c\vec{x}\|_1=\sum_{i=1}^{N}|cx_i|=\sum_{i=1}^{N}|c|\,|x_i|=|c|\sum_{i=1}^{N}|x_i|=|c|\,\|\vec{x}\|_1.
\]
This satisfies nonnegative homogeneity.
\item Using the triangle inequality $|a+b|\le |a|+|b|$,
\[
\|\vec{x}+\vec{y}\|_1=\sum_{i=1}^{N}|x_i+y_i|\le \sum_{i=1}^{N}\big(|x_i|+|y_i|\big)=\sum_{i=1}^{N}|x_i|+\sum_{i=1}^{N}|y_i|=\|\vec{x}\|_1+\|\vec{y}\|_1.
\]
This satisfies the triangle inequality.
\item Since each $|x_i|\ge 0$, we have $\|\vec{x}\|_1=\sum_{i=1}^{N}|x_i|\ge 0$. Thus, (I think fairly trivially),
\[
\|\vec{x}\|_1=0
\iff
\sum_{i=1}^{N}|x_i|=0
\iff
|x_i|=0\ \text{for all }i
\iff
x_i=0\ \text{for all }i
\iff
\vec{x}=\vec{0}.
\]
This satisfies nonnegative definiteness.
\end{enumerate}
Therefore, the one-norm is a vector norm.
\item Similarly, let us check the three axioms for the infinity-norm.
\begin{enumerate}[label=(\roman*)]
\item For any scalar $c\in\mathbb{R}$,
\[
\|c\vec{x}\|_\infty=\max\{|cx_1|,\dots,|cx_N|\}
=\max\{|c||x_1|,\dots,|c||x_N|\}
=|c|\max\{|x_1|,\dots,|x_N|\}
=|c|\,\|\vec{x}\|_\infty.
\]
\item For each $i$, we use $|a+b|\le |a|+|b|$ to get
\[
|x_i+y_i|\le |x_i|+|y_i|.
\]
Also, $\|\vec{x}\|_\infty=\max\{|x_1|,\dots,|x_N|\}$, so $|x_i|\le \|\vec{x}\|_\infty$ for every $i$, and similarly $|y_i|\le \|\vec{y}\|_\infty$ for every $i$. Thus, for every $i$,
\[
|x_i+y_i|\le |x_i|+|y_i|\le \|\vec{x}\|_\infty+\|\vec{y}\|_\infty.
\]
Since this inequality holds for every $i$, the maximum of the left-hand side over all $i$ is also at most $\|\vec{x}\|_\infty+\|\vec{y}\|_\infty$, so
\[
\|\vec{x}+\vec{y}\|_\infty=\max\{|x_1+y_1|,\dots,|x_N+y_N|\}\le \|\vec{x}\|_\infty+\|\vec{y}\|_\infty.
\]
\item Since each $|x_i|\ge 0$, it follows that $\|\vec{x}\|_\infty=\max\{|x_1|,\dots,|x_N|\}\ge 0$. Thus,
\[
\|\vec{x}\|_\infty=0
\iff
\max\{|x_1|,\dots,|x_N|\}=0
\iff
|x_i|=0\ \text{for all }i
\iff
x_i=0\ \text{for all }i
\iff
\vec{x}=\vec{0}.
\]
\end{enumerate}
Therefore, the infinity-norm is a vector norm.\end{enumerate}

\begin{problem}
[Problem 5]
Suppose $\vec{x}\in\mathbb{R}^{100}$ with $\mathrm{rms}(\vec{x})=1$.
\begin{enumerate}[label=(\alph*)]
\item What is the maximum number of entries of $\vec{x}$ that can satisfy $|x_i|\ge 3$? Call this quantity $k$. Explain why no such vector can have $k+1$ entries with absolute value at least $3$.
\item Show in this case that Chebyshev's inequality is sharp, i.e. give an example of a specific $\vec{x}\in\mathbb{R}^{100}$ with RMS value $1$ and $k$ entries larger than $3$ in absolute value.
\end{enumerate}
\end{problem}

\begin{enumerate}[label=(\alph*)]
\item Let $k$ be the number of indices $i\in\{1,\dots,100\}$ such that $|x_i|\ge 3$.

Since $\mathrm{rms}(\vec{x})=1$ and $N=100$, by definition
\[
\mathrm{rms}(\vec{x})=\frac{\|\vec{x}\|}{\sqrt{100}}
\qquad\Rightarrow\qquad
1=\frac{\|\vec{x}\|}{10}
\qquad\Rightarrow\qquad
\|\vec{x}\|=10
\qquad\Rightarrow\qquad
\|\vec{x}\|^2=100.
\]

Since Chebyshev's Inequality states for any $\alpha>0$, the number of indices with $|x_i|\ge \alpha$ is at most
\[
\frac{\|\vec{x}\|^2}{\alpha^2},
\]
applying this with $\alpha=3$, we get
\[
k \le \frac{\|\vec{x}\|^2}{3^2}=\frac{100}{9}.
\]
Since $k$ is an integer, the largest possible value is
\[
k=\left\lfloor \frac{100}{9}\right\rfloor=11.
\]
Finally, no such vector can have $k+1=12$ entries with $|x_i|\ge 3$, because then
\[
\|\vec{x}\|^2=\sum_{i=1}^{100}x_i^2 \ge 12\cdot 3^2=108>100.
\]
\item Consider the vector $\vec{x}\in\mathbb{R}^{100}$ where
\[
x_1=\cdots=x_{11}=3,\qquad x_{12}=1,\qquad x_{13}=\cdots=x_{100}=0.
\]
Then, exactly $k=11$ entries satisfy $|x_i|\ge 3$. Also,
\[
\|\vec{x}\|^2=\sum_{i=1}^{100}x_i^2=11\cdot 3^2+1^2=99+1=100,
\]
so $\|\vec{x}\|=10$, and thus,
\[
\mathrm{rms}(\vec{x})=\frac{\|\vec{x}\|}{\sqrt{100}}=\frac{10}{10}=1.
\]
Therefore, Chebyshev's Inequality is sharp.
\end{enumerate}

\begin{problem}
[Problem 6]
Let $\vec{x}\in\mathbb{R}^N$ and $\alpha,\beta\in\mathbb{R}$ be scalars. Show that
\begin{enumerate}[label=(\alph*)]
\item $\mathrm{avg}(\alpha\vec{x}+\beta\vec{\mathbf{1}})=\alpha\,\mathrm{avg}(\vec{x})+\beta$, and
\item $\mathrm{std}(\alpha\vec{x}+\beta\vec{\mathbf{1}})=|\alpha|\,\mathrm{std}(\vec{x})$.
\end{enumerate}
\end{problem}

\begin{enumerate}[label=(\alph*)]
\item By definition,
\[
\mathrm{avg}(\vec{x})=\frac{1}{N}\vec{\mathbf{1}}\cdot\vec{x}.
\]
Thus,
\[
\mathrm{avg}(\alpha\vec{x}+\beta\vec{\mathbf{1}})
=\frac{1}{N}\vec{\mathbf{1}}\cdot(\alpha\vec{x}+\beta\vec{\mathbf{1}})
=\frac{1}{N}\big(\alpha(\vec{\mathbf{1}}\cdot\vec{x})+\beta(\vec{\mathbf{1}}\cdot\vec{\mathbf{1}})\big).
\]
Since $\vec{\mathbf{1}}=\begin{bmatrix}1\\ \vdots\\ 1\end{bmatrix}\in\mathbb{R}^N$,
\[
\vec{\mathbf{1}}\cdot\vec{\mathbf{1}}=1^2+\cdots+1^2=N,
\]
it follows that
\begin{align*}
\mathrm{avg}(\alpha\vec{x}+\beta\vec{\mathbf{1}})
&=\frac{1}{N}\big(\alpha(\vec{\mathbf{1}}\cdot\vec{x})+\beta(\vec{\mathbf{1}}\cdot\vec{\mathbf{1}})\big)\\
&=\frac{1}{N}\big(\alpha(\vec{\mathbf{1}}\cdot\vec{x})+\beta N\big)\\
&=\alpha\frac{1}{N}(\vec{\mathbf{1}}\cdot\vec{x})+\beta\\
&=\alpha\,\mathrm{avg}(\vec{x})+\beta.
\end{align*}
\item By definition,
\[
\mathrm{std}(\vec{x})=\mathrm{rms}\big(\vec{x}-\mathrm{avg}(\vec{x})\vec{\mathbf{1}}\big)
=\frac{1}{\sqrt{N}}\big\|\vec{x}-\mathrm{avg}(\vec{x})\vec{\mathbf{1}}\big\|.
\]
Thus,
\[
\mathrm{std}(\alpha\vec{x}+\beta\vec{\mathbf{1}})
=\frac{1}{\sqrt{N}}\left\|(\alpha\vec{x}+\beta\vec{\mathbf{1}})-\mathrm{avg}(\alpha\vec{x}+\beta\vec{\mathbf{1}})\vec{\mathbf{1}}\right\|.
\]
From part (a),
\[
\mathrm{avg}(\alpha\vec{x}+\beta\vec{\mathbf{1}})=\alpha\,\mathrm{avg}(\vec{x})+\beta,
\]
so
\begin{align*}
(\alpha\vec{x}+\beta\vec{\mathbf{1}})-\mathrm{avg}(\alpha\vec{x}+\beta\vec{\mathbf{1}})\vec{\mathbf{1}}
&=(\alpha\vec{x}+\beta\vec{\mathbf{1}})-(\alpha\,\mathrm{avg}(\vec{x})+\beta)\vec{\mathbf{1}}\\
&=(\alpha\vec{x}+\beta\vec{\mathbf{1}})-\big(\alpha\,\mathrm{avg}(\vec{x})\vec{\mathbf{1}}+\beta\vec{\mathbf{1}}\big)\\
&=\alpha\vec{x}+\beta\vec{\mathbf{1}}-\alpha\,\mathrm{avg}(\vec{x})\vec{\mathbf{1}}-\beta\vec{\mathbf{1}}\\
&=\alpha\big(\vec{x}-\mathrm{avg}(\vec{x})\vec{\mathbf{1}}\big).
\end{align*}
Therefore,
\begin{align*}
\mathrm{std}(\alpha\vec{x}+\beta\vec{\mathbf{1}})
&=\frac{1}{\sqrt{N}}\left\|\alpha\big(\vec{x}-\mathrm{avg}(\vec{x})\vec{\mathbf{1}}\big)\right\|\\
&=\frac{1}{\sqrt{N}}|\alpha|\,\big\|\vec{x}-\mathrm{avg}(\vec{x})\vec{\mathbf{1}}\big\|\\
&=|\alpha|\,\mathrm{std}(\vec{x}).
\end{align*}
\end{enumerate}
\end{document}